\chapter{Marchés d'actions asiatiques et émergents}\label{ch:m1:asian-and-emerging-equity-markets}

À Shanghai, une action qui a gagné dix pour cent depuis la clôture de la veille
ne peut plus monter aujourd'hui, et celui qui l'a achetée ce matin ne peut la revendre
que demain. À Hong Kong, l'acheteur et le vendeur versent chacun 0.1~\% de chaque transaction à l'État.
À Tokyo, le plus petit pas de cotation d'une action
dépend de l'indice auquel elle appartient. À Mumbai, le vendeur d'une option paie
une taxe sur la prime, une transaction peut être réglée dans l'après-midi même, et un
fonds étranger peut constater que les étrangers, pris ensemble, détiennent déjà autant
de la société que la loi l'autorise. Une stratégie qui fonctionne à New York ne peut pas
être «~transposée~» telle quelle vers ces marchés~; elle doit être redérivée selon un autre
ensemble de règles. Ce chapitre est un guide de terrain des règles qui diffèrent, et une méthode
pour lire tout nouveau marché.

\section{Ce qui peut différer}

\begin{method}[Lire un nouveau marché d'actions]\label{met:m1:asian-and-emerging-equity-markets:checklist}
Avant tout chargement de données, remplir ce tableau pour le marché considéré.
\begin{enumerate}
\item \textbf{Accès~:} qui peut détenir directement des actions~; licences, quotas,
  plafonds de détention~; si un investisseur étranger négocie la même ligne de titres
  que celle négociée par un investisseur local.
\item \textbf{Unité et pas de cotation~:} le \omterm{def:m1:asian-and-emerging-equity-markets:lot}{lot de négociation} dans lequel les ordres doivent être exprimés~; la
  table des pas de cotation.
\item \textbf{Limites et suspensions~:} limites quotidiennes de variation de cours, coupe-circuits, leur
  cours de référence.
\item \textbf{Revente~:} si les actions achetées aujourd'hui peuvent être revendues le jour même~;
  si, et comment, elles peuvent être vendues à découvert.
\item \textbf{Règlement-livraison et trésorerie~:} cycle, préfinancement, convertibilité
  de la devise.
\item \textbf{Impôts et frais~:} taxes par transaction selon l'instrument et le sens~;
  retenue à la source sur les dividendes~; gains en capital.
\item \textbf{Séances~:} enchères, pauses méridiennes, séances à prix fixe
  d'après-clôture.
\end{enumerate}
Chaque ligne modifie l'éventail des stratégies possibles~; plusieurs d'entre elles modifient
les propriétés statistiques des données elles-mêmes.
\end{method}

\section{Limites de variation de cours}

\begin{definition}[Limite de variation de cours]\label{def:m1:asian-and-emerging-equity-markets:limit}
Une \emph{limite de variation de cours}\index{limite de variation de cours} interdit les ordres et les transactions
au-delà d'un écart fixe, en pourcentage ou en montant, par rapport à un cours de référence,
généralement la clôture de la veille. Une action à sa limite haute est \emph{en limite haute}~:
les ordres d'achat se mettent en file au cours limite et ne s'exécutent que si quelqu'un vend à ce prix.
\end{definition}

\begin{proposition}[Jours bloqués à la limite]\label{prop:m1:asian-and-emerging-equity-markets:locked}
Une nouvelle modifie la fair value d'une action d'un facteur $1+S$ avec $S>0$. Sous une
limite quotidienne $L$, l'action passe
\[
n = \left\lceil \frac{\ln(1+S)}{\ln(1+L)} \right\rceil - 1
\]
jours entiers où elle est bloquée à la limite haute avant un jour où elle peut s'échanger
à sa fair value. Pour une baisse, on remplace $1+S$ et $1+L$ par $1-S$ et $1-L$.
\end{proposition}
\begin{proof}
Après $k$ jours à la limite, le cours vaut au plus $(1+L)^k$ fois le cours initial.
La fair value devient atteignable au premier jour $k$ tel que $(1+L)^k \ge 1+S$~;
les $k-1$ jours précédents sont bloqués.
\end{proof}

\begin{example}[Les baisses prennent plus de temps]\label{ex:m1:asian-and-emerging-equity-markets:asym}
Avec $L = 10\%$, une hausse de 50~\% nécessite $\lceil 4.25\rceil = 5$ jours, dont quatre
bloqués. Une baisse de 50~\% nécessite $\lceil \ln 0.5/\ln 0.9\rceil = \lceil
6.58\rceil = 7$ jours, dont six bloqués~: pendant six jours, les détenteurs souhaitant vendre
ne le peuvent pas, à aucun prix. Un fonds qui promet à ses propres investisseurs une liquidité quotidienne
et qui détient de tels titres pose un problème qu'aucun modèle de risque construit sur des
rendements quotidiens n'aura mesuré (voir la \cref{fig:m1:asian-and-emerging-equity-markets:locked}).
\end{example}

\begin{omfigure}
\begin{tikzpicture}
\begin{axis}[width=10.5cm,height=4.6cm,
  xlabel={ampleur de la réévaluation à la hausse (\%)},ylabel={jours bloqués en limite haute},
  tick label style={font=\small},label style={font=\small},
  xmin=0,xmax=80,ymin=0,ymax=12.5,grid=major,grid style={gray!25},
  legend columns=3,legend style={font=\footnotesize,at={(0.5,-0.36)},anchor=north,draw=none,fill=none,/tikz/every even column/.append style={column sep=8pt}},
  table/col sep=comma]
\addplot[omThm,very thick,const plot] table[x=shock_pct,y=l5]{figdata/markets-1/12-asian-and-emerging-equity-markets/locked.csv};
\addplot[omDef,very thick,const plot] table[x=shock_pct,y=l10]{figdata/markets-1/12-asian-and-emerging-equity-markets/locked.csv};
\addplot[omMeth,very thick,const plot] table[x=shock_pct,y=l20]{figdata/markets-1/12-asian-and-emerging-equity-markets/locked.csv};
\legend{limite 5~\%,limite 10~\%,limite 20~\%}
\end{axis}
\end{tikzpicture}

\omcaption{Jours complets passés à la limite avant qu'une réévaluation puisse être négociée,
selon la \cref{prop:m1:asian-and-emerging-equity-markets:locked}. Données~: calculées par
le script du chapitre.}\label{fig:m1:asian-and-emerging-equity-markets:locked}
\end{omfigure}

\begin{proposition}[Effet des limites sur les données]\label{prop:m1:asian-and-emerging-equity-markets:data}
Soient des rendements quotidiens réels indépendants. Les rendements observés sous une \omterm{def:m1:asian-and-emerging-equity-markets:limit}{limite
de variation de cours} (i) ont une variance plus faible que celle des rendements réels, (ii) sont
positivement autocorrélés, et (iii) présentent des masses ponctuelles en $\pm L$.
\end{proposition}
\begin{proof}[Preuve partielle]
(iii) découle de la définition. Pour (i) et (ii)~: un jour à la limite, le rendement observé
vaut $L$ et la part non exprimée du mouvement est reportée au jour suivant, où
elle s'ajoute au rendement indépendant de ce jour-là. La part reportée a le signe du
rendement du jour, de sorte que des rendements observés consécutifs covarient positivement~; et comme
un grand mouvement est réparti sur plusieurs jours, la somme des carrés des rendements
quotidiens est inférieure au carré du mouvement. Un calcul complet exige la
distribution des rendements, et il est laissé à la simulation du tutoriel.
\end{proof}

\begin{omfigure}
\begin{tikzpicture}
\begin{axis}[width=11cm,height=5cm,ybar=0pt,bar width=3.4pt,ymode=log,log origin=infty,
  xlabel={rendement quotidien (\%)},ylabel={jours sur 60\,000 (échelle log)},
  tick label style={font=\small},label style={font=\small},
  xmin=-14,xmax=14,ymin=1,ymax=30000,
  legend columns=2,legend style={font=\footnotesize,at={(0.5,-0.32)},anchor=north,draw=none,fill=none,/tikz/every even column/.append style={column sep=8pt}},
  area legend,table/col sep=comma]
\addplot[fill=omExo!50,draw=omExo] table[x=centre_pct,y=true]{figdata/markets-1/12-asian-and-emerging-equity-markets/limit_hist.csv};
\addplot[fill=omDef!70,draw=omDef] table[x=centre_pct,y=observed]{figdata/markets-1/12-asian-and-emerging-equity-markets/limit_hist.csv};
\legend{rendements réels,observés sous une limite de 10~\%}
\end{axis}
\end{tikzpicture}

\omcaption{Rendements à queue épaisse, avec et sans limite quotidienne de 10~\%. La
série observée n'a aucune queue au-delà de la limite et présente un pic à celle-ci~; son écart type
vaut 3.03~\% contre 3.45~\%, et son autocorrélation d'ordre 1, $+0.09$, contre zéro.
Données~: la simulation du tutoriel.}\label{fig:m1:asian-and-emerging-equity-markets:hist}
\end{omfigure}

Une volatilité estimée sur des clôtures contraintes par la limite est trop faible~; un signal de momentum
ajusté sur ces données repère une «~continuation~» qu'il est impossible d'exploiter, car
les jours où il la prédit, personne ne peut acheter. Ces deux erreurs figurent parmi les
plus courantes de la recherche intermarchés.

\section{Chine continentale et le Connect}

\begin{definition}[Restriction T+0]\label{def:m1:asian-and-emerging-equity-markets:tplus0}
Une \emph{restriction T+0}\index{restriction T+0} (en Chine continentale, la
«~règle de négociation T+1~») interdit de vendre des actions le jour même de leur achat. Il
s'agit d'une règle de négociation, distincte du \omterm{def:m1:clearing-and-settlement:cycle}{cycle de règlement-livraison}.
\end{definition}

\begin{definition}[Investisseur étranger qualifié et Stock Connect]\label{def:m1:asian-and-emerging-equity-markets:connect}
Un \emph{investisseur étranger qualifié}\index{investisseur étranger qualifié} est une
institution étrangère agréée par le régulateur de Chine continentale pour investir directement dans les
titres cotés sur le marché intérieur. \emph{Stock Connect}\index{Stock Connect} est le lien entre la
bourse de Hong Kong et les bourses de Shanghai et de Shenzhen, qui permet aux investisseurs
de chaque côté de négocier les actions éligibles de l'autre marché par l'intermédiaire de leurs
courtiers et chambres de compensation d'origine, dans la limite d'un quota quotidien, sans agrément.
\end{definition}

\begin{dated}{2026-09}{Règles de négociation des actions A continentales}\label{dat:m1:asian-and-emerging-equity-markets:china}
Les actions du marché principal de Shanghai et de Shenzhen peuvent varier de 10~\% par jour~; celles
du STAR Market et du ChiNext de 20~\%, sans limite pendant les cinq premiers jours de cotation
d'une valeur au STAR. Les règles de négociation révisées des bourses de Shanghai, de Shenzhen et de
Pékin sont entrées en vigueur le 6 juillet 2026~: la limite applicable aux actions du marché principal
sous régime spécial a été portée à 10~\%, et la séance à prix fixe après la clôture (de 15:05 à
15:30) a été étendue à toutes les actions A et à tous les fonds indiciels cotés (ETF).
Les actions achetées ne peuvent pas être vendues le jour même. Le quota \omterm{def:m1:asian-and-emerging-equity-markets:connect}{Stock Connect} en sens nord
est de 52~milliards de renminbis par jour pour chacun des liens avec Shanghai et Shenzhen.
\end{dated}

\begin{omfigure}
\begin{tikzpicture}[font=\small,
  box/.style={draw=omDef,rounded corners=2pt,minimum width=3.0cm,minimum height=1.0cm,align=center,fill=omDef!4},
  mid/.style={draw=omThm,rounded corners=2pt,minimum width=3.0cm,minimum height=1.0cm,align=center,fill=omThm!5},
  flow/.style={-{Stealth[length=2mm]},thick}]
\node[box] (f) at (0,0) {Investisseur\\international};
\node[box] (hb) at (4.2,0) {Courtier\\de Hong Kong};
\node[mid] (hk) at (8.4,0) {Bourse de Hong Kong\\et chambre de compensation};
\node[box] (sh) at (12.8,0.9) {Bourse de\\Shanghai};
\node[box] (sz) at (12.8,-0.9) {Bourse de\\Shenzhen};
\draw[flow,omDef] (f) -- (hb);
\draw[flow,omDef] (hb) -- (hk);
\draw[flow,omDef] (hk) -- (sh);
\draw[flow,omDef] (hk) -- (sz);
\node[font=\footnotesize,omThm,align=center] at (8.4,-1.25) {quota quotidien contrôlé ici~;\\règlement en renminbis};
\end{tikzpicture}

\omcaption{Un ordre \omterm{def:m1:asian-and-emerging-equity-markets:connect}{Stock Connect} en sens nord. L'investisseur n'a pour interlocuteur que son
courtier de Hong Kong~; la filiale de la bourse de Hong Kong est membre de la bourse continentale,
et les deux chambres de compensation sont contreparties l'une de l'autre. Les règles continentales
--- limites, lots, interdiction de vendre le jour même
--- s'appliquent à l'ordre dans son intégralité.}
\end{omfigure}

\begin{remark}[La même société à deux cours]\label{rem:m1:asian-and-emerging-equity-markets:ah}
Un grand nombre de sociétés continentales sont cotées à la fois sous forme d'actions A sur le marché intérieur
et d'actions H à Hong Kong. Les deux lignes donnent droit aux mêmes dividendes et aux mêmes droits de vote et
ne sont \emph{pas} convertibles l'une en l'autre~; elles se négocient couramment à des cours différents, car
les catégories d'investisseurs autorisées à détenir chacune d'elles sont différentes et aucun arbitragiste
ne peut livrer l'une contre l'autre. Un écart de cours qui ne peut pas être résorbé par la livraison
relève de l'accès au marché, et non d'un gain sans risque
(\cref{ch:m1:delta-one-instruments}).
\end{remark}

\section{Hong Kong}

\begin{definition}[Lot de négociation]\label{def:m1:asian-and-emerging-equity-markets:lot}
Un \emph{lot de négociation}\index{lot de négociation} est le nombre d'actions dont les ordres
du carnet d'ordres principal d'une bourse doivent être des multiples. Lorsque l'émetteur
le choisit, le lot de négociation varie d'une action à l'autre.
\end{definition}

\begin{definition}[Droit de timbre]\label{def:m1:asian-and-emerging-equity-markets:stamp}
Un \emph{droit de timbre}\index{droit de timbre} sur les transferts d'actions est un impôt prélevé en
pourcentage de la valeur de chaque transaction, à la charge de l'acheteur, du vendeur ou des deux,
quel que soit le gain ou la perte.
\end{definition}

\begin{dated}{2026-09}{Hong Kong~: droit et lots}\label{dat:m1:asian-and-emerging-equity-markets:hk}
Le \omterm{def:m1:asian-and-emerging-equity-markets:stamp}{droit de timbre} sur les transferts d'actions à Hong Kong est de 0.1~\% de la valeur pour l'acheteur
et de 0.1~\% pour le vendeur, depuis le 17 novembre 2023 (contre 0.13~\% auparavant). Les \omterm{def:m1:asian-and-emerging-equity-markets:lot}{lots de
négociation} étaient choisis librement par les émetteurs, en plus de quarante tailles différentes~; une
réforme mise en œuvre à partir de juillet 2026 les limite à huit tailles standard (1, 50, 100, 500,
1\,000, 2\,000, 5\,000 et 10\,000 actions), avec un minimum plus bas et une nouvelle valeur maximale
par lot.
\end{dated}

\begin{example}[Vingt points de base]\label{ex:m1:asian-and-emerging-equity-markets:duty}
Un aller-retour sur une action de Hong Kong coûte 20 points de base de droits avant toute
fourchette, commission ou impact. Une stratégie de \omterm{def:m1:what-a-trading-firm-does:market-maker}{teneur de marché} qui capte un ou deux points
de base par aller-retour
(\cref{dat:m1:what-a-trading-firm-does:listed}) ne peut pas exister sous cet impôt~;
des stratégies à horizon de plusieurs jours, qui gagnent cinquante points de base par position, le peuvent.
Les taxes sur les transactions ne pèsent pas de la même façon sur toutes les stratégies~: elles
suppriment l'extrémité à court horizon du spectre, et avec elle les acteurs qui auraient fourni la
liquidité intrajournalière.
\end{example}

\section{Japon, Corée et la règle de l'uptick}

La Bourse de Tokyo négocie les actions par unités de 100 et publie plusieurs tables
de tick~; les valeurs composant ses indices à forte capitalisation utilisent une table plus fine
que celle des autres actions, jusqu'à une fraction de yen pour les plus liquides. Une action qui
entre dans l'indice ou en sort change de table, ce qui modifie sa fourchette, la longueur de ses
files d'attente et la rentabilité de la tenue de marché sur ce titre, sans aucun changement
dans la société.

\begin{definition}[Règle de l'uptick]\label{def:m1:asian-and-emerging-equity-markets:uptick}
Une \emph{règle de l'uptick}\index{règle de l'uptick} (uptick rule) restreint les prix auxquels une vente à découvert
peut être passée~: en général pas en dessous du dernier cours échangé, ou pas en dessous de ce cours
une fois que l'action a baissé d'un montant donné au cours de la journée.
\end{definition}

\begin{example}[Dix-sept mois sans vente à découvert]\label{ex:m1:asian-and-emerging-equity-markets:korea}
En novembre 2023, les autorités coréennes ont interdit la vente à découvert de toutes les actions cotées,
en invoquant les ventes à découvert à nu des banques mondiales. L'interdiction a pris fin le 31~mars
2025, date à laquelle la vente à découvert a été rétablie pour toutes les actions cotées, en même temps
qu'un nouveau système électronique de détection des ventes à découvert à nu. Pendant dix-sept mois,
aucune stratégie long--short sur actions n'a pu être menée sur le territoire d'un des plus grands marchés mondiaux~;
l'arbitrage sur indices, la couverture des obligations convertibles et le market making sur les dérivés sur actions ont été
entravés simultanément. Un risque réglementaire de ce type est un paramètre du marché, au même
titre que la volatilité.
\end{example}

\section{Inde}

\begin{definition}[Taxe sur les transactions de titres et plafond de détention étrangère]\label{def:m1:asian-and-emerging-equity-markets:stt}
Une \emph{taxe sur les transactions de titres}\index{taxe sur les transactions de titres} (securities transaction tax) (STT)
est un impôt sur les transactions en bourse dont le taux et le redevable dépendent de
l'instrument~: actions, contrats à terme ou options. Un
\emph{plafond de détention étrangère}\index{plafond de détention étrangère} (foreign ownership limit) limite la
fraction d'une société que des non-résidents peuvent détenir au total~; une fois ce plafond atteint, les
étrangers ne peuvent acheter qu'à d'autres étrangers.
\end{definition}

\begin{dated}{2026-09}{Inde~: fiscalité et règlement-livraison}\label{dat:m1:asian-and-emerging-equity-markets:india}
La STT sur les actions achetées ou vendues pour livraison est de 0.1~\% de la valeur de chaque côté.
À partir du 1\up{er}~avril 2026, le taux sur les contrats à terme est de 0.05~\% de la valeur, à la charge
du vendeur (précédemment 0.02~\%), et sur les options de 0.15~\% de la prime, à la charge du
vendeur (précédemment 0.10~\%). Les actions se règlent à $T{+}1$~; un cycle optionnel le jour même
($T{+}0$) est disponible pour les 500 plus grandes actions.
\end{dated}

\begin{omfigure}
\begin{tikzpicture}
\begin{axis}[width=9cm,height=3.6cm,xbar,bar width=8pt,
  xlabel={taxe sur une transaction aller-retour (points de base de la valeur)},
  ytick=data,yticklabels={{Actions de Hong Kong (droit de timbre)},{Actions indiennes livrées (STT)},{Contrats à terme sur indice indiens (STT)}},
  yticklabel style={font=\small},tick label style={font=\small},label style={font=\small},
  y dir=reverse,xmin=0,xmax=25,enlarge y limits=0.25,grid=major,grid style={gray!25},
  nodes near coords,nodes near coords style={font=\footnotesize},
  table/col sep=comma]
\addplot[fill=omDef!60,draw=omDef] table[x=bp,y=k]{figdata/markets-1/12-asian-and-emerging-equity-markets/taxes.csv};
\end{axis}
\end{tikzpicture}

\omcaption{Taxes par transaction sur un achat suivi d'une vente. Un marché sans
une telle taxe afficherait zéro. Données~: les taux figurant dans
les \cref{dat:m1:asian-and-emerging-equity-markets:hk,dat:m1:asian-and-emerging-equity-markets:india}.}
\end{omfigure}

\section{Ce qui change pour une société de trading}

\begin{proposition}[Le plancher fiscal sur les volumes échangés]\label{prop:m1:asian-and-emerging-equity-markets:floor}
Une stratégie dégage un avantage brut $e$ par aller-retour et supporte une
taxe de transaction aller-retour $\tau$ ainsi que d'autres coûts $\kappa$, le tout exprimé en fractions de la valeur.
Elle n'est viable que si $e > \tau + \kappa$~; et si son avantage par unité de durée de détention
est à peu près constant, $e \approx \alpha h$ pour une durée de détention $h$, la
plus courte durée de détention viable est $h^\star = (\tau+\kappa)/\alpha$.
\end{proposition}
\begin{proof}
Le profit net par aller-retour est $e - \tau - \kappa$~; on remplace $e = \alpha h$.
\end{proof}

Avec un $\alpha$ de 5 points de base par jour et un $\kappa$ de 5 points de base, un
marché sans taxe admet des durées de détention à partir d'un jour~; un marché avec 20
points de base de droits ne les admet qu'à partir de cinq jours. Le même signal qui est négocié
en cours de journée sur un marché doit être ralenti, voire abandonné, sur un autre.

\section{Tutoriel~: effet d'une limite de variation de cours sur les rendements}\label{tut:m1:asian-and-emerging-equity-markets:limits}

\begin{tutorial}
\textbf{Objectif.} Faire passer une série de rendements vrais à travers une limite de 10~\% et mesurer
ce qu'elle fait à la volatilité et à l'autocorrélation. \textbf{État final~:}
la \cref{fig:m1:asian-and-emerging-equity-markets:hist} et ses trois
statistiques.
\begin{tutsteps}
\item \textbf{La limite avec report.} Le cours observé ne peut pas varier de plus de $L$ par jour~;
le reste du mouvement attend.
\omcode{code/markets-1/12-asian-and-emerging-equity-markets/python/asia_rules.py}{17}{27}{Troncature avec report, en espace logarithmique afin que le mouvement total soit conservé.}\label{lst:m1:asian-and-emerging-equity-markets:truncate}
\item \textbf{Vérifier la conservation.} Une seule réévaluation de $+35\%$ devient
$+10\%$, $+10\%$, $+10\%$ et un reliquat~; le produit des rendements bruts observés
vaut 1.35.
\item \textbf{Mesurer.} Avec des rendements à queues lourdes d'écart type quotidien de 3.45~\%,
2~\% des jours se terminent à une limite~; l'écart type observé vaut 3.03~\% et la première
autocorrélation $+0.088$.
\item \textbf{Jours bloqués.}
\omcode{code/markets-1/12-asian-and-emerging-equity-markets/python/asia_rules.py}{7}{14}{La proposition sous forme de fonction, avec l'asymétrie entre hausses et baisses.}\label{lst:m1:asian-and-emerging-equity-markets:locked}
\end{tutsteps}
\textbf{Pour aller plus loin.} Ajuster une règle de momentum à un jour sur la série observée et
calculer son «~profit~»~; interdire ensuite les transactions les jours de limite et recalculer.
Élargir ensuite la limite à 20~\% et voir quelle
part de la distorsion subsiste.
\end{tutorial}

\section{Construction~: la table des règles de marché}\label{bld:m1:asian-and-emerging-equity-markets:rules}

\begin{build}
\bfield{Rôle} Chaque ordre créé par la petite société de trading est validé, avant son envoi,
par rapport aux règles du marché vers lequel il se dirige~; chaque backtest interroge la même table pour
savoir ce qui était autorisé à cette date. \bfield{Interface} \texttt{MarketRules(market, valid\_from, lot, tick\_table,
limit\_up, limit\_down, same\_day\_sell, short\_selling, settlement\_days, tax\_buy, tax\_sell)}~; \texttt{RuleBook.load(csv)}~; \texttt{rules(market, date)}~; \texttt{validate(order, reference\_price,
bought\_today)} renvoyant une liste d'infractions~; \texttt{tax(side, value)}.
\bfield{Règles} Les règles sont versionnées par date~: une requête pour une date passée renvoie les
règles alors en vigueur, jamais celles d'aujourd'hui. Les quantités doivent être des multiples du
lot~; les prix, des multiples du tick applicable à ce prix et situés à l'intérieur de la bande de
limites autour du cours de référence~; une vente excédant la position détenue au début de la journée
est rejetée là où la vente le jour même est interdite.
\bfield{Tests d'acceptation} \texttt{code/firm/marketrules/tests/}~: arrondi au lot et au tick~; bandes
de limites~; rejet de la vente le jour même~; un changement de taux à une date donnée n'est appliqué
qu'à partir de cette date.
\bfield{Pour aller plus loin} Charger l'historique des règles d'un marché réel à partir des circulaires de sa bourse et
reproduire les encadrés datés de ce chapitre sous forme de requêtes.
\end{build}

\begin{omsources}
\item KPMG, \emph{Hong Kong: Reduced stamp duty rate on stock transfers
effective 17 November 2023}~; HKEX, consultation et conclusions sur le
cadre du lot de négociation, 2025--2026.
\item HKEX, \emph{Stock Connect: Information Booklet and FAQ}, 2024.
\item Bourses de Shanghai, de Shenzhen et de Pékin, règles de négociation
révisées, en vigueur le 6 juillet 2026 (telles que rapportées dans la presse financière).
\item Financial Services Commission of Korea, \emph{Stock Short Selling to be
Fully Reinstated from March 31}, communiqué de presse, 2025.
\item ClearTax, \emph{Securities Transaction Tax: latest updates}, et ICICI
Direct, \emph{STT changes in Budget 2026} (taux applicables à partir du 1\up{er}~avril 2026)~;
Securities and Exchange Board of India, circulaires sur le cycle de
règlement-livraison $T{+}0$ optionnel.
\item Japan Exchange Group, \emph{Trading Rules of Domestic Stocks: Tick
Size}.
\item K.~Kim et S.~Rhee, «~Price limit performance: evidence from the Tokyo
Stock Exchange~», \emph{Journal of Finance} 52 (1997).
\end{omsources}

\section{Exercices}

\begin{exercise}[$\star$]\label{exo:m1:asian-and-emerging-equity-markets:1}
Une action a clôturé à 24.00 avec une limite de 10~\% et un tick de 0.01.
Donner les cours plafond et plancher d'aujourd'hui. Le système d'un courtier calcule le plafond
à 26.40 exactement~; selon quelle convention d'arrondi celui de la bourse pourrait-il différer~?
different~?
\end{exercise}

\begin{exercise}[$\star$]\label{exo:m1:asian-and-emerging-equity-markets:2}
Un investisseur souhaite acquérir pour 250\,000~HKD d'une action qui se négocie à 61.30~HKD
avec un \omterm{def:m1:asian-and-emerging-equity-markets:lot}{lot de négociation} de 500. Calculer le nombre d'actions qu'il peut acheter, leur valeur
et le \omterm{def:m1:asian-and-emerging-equity-markets:stamp}{droit de timbre} qu'il paie à l'achat.
\end{exercise}

\begin{exercise}[$\star$]\label{exo:m1:asian-and-emerging-equity-markets:3}
À l'aide de l'\cref{dat:m1:asian-and-emerging-equity-markets:india}, calculer l'impôt
sur (a) l'achat puis la revente ultérieure, avec livraison, d'actions pour 5~millions de roupies indiennes~; (b)
la vente de contrats à terme sur indice d'une valeur contractuelle de 20~millions de roupies indiennes~; (c) la vente
d'options pour une prime de 400\,000~roupies indiennes.
\end{exercise}

\begin{exercise}[$\star\star$]\label{exo:m1:asian-and-emerging-equity-markets:4}
Une mauvaise nouvelle réduit de 35~\% la fair value d'une action. Calculer pendant
combien de jours elle reste cotée en limite basse avec une limite de 10~\%, de 20~\% puis de 5~\%.
\end{exercise}

\begin{exercise}[$\star\star$]\label{exo:m1:asian-and-emerging-equity-markets:5}
Un signal rapporte un avantage brut de 4~points de base par jour de détention~; les coûts autres que
la taxe s'élèvent à 6~points de base par aller-retour.
Calculer la plus courte durée de détention viable sans taxe sur les transactions, puis avec un droit de 20~points de base par aller-retour, et
le profit net par aller-retour pour une détention de dix jours dans chaque cas.
\end{exercise}

\begin{exercise}[$\star\star$]\label{exo:m1:asian-and-emerging-equity-markets:6}
Une trader opérant sur un marché où la vente le jour même est interdite achète 10\,000 actions
à l'ouverture~; à midi, l'action progresse de 6~\% et elle veut sécuriser le gain.
L'action dispose d'un contrat à terme liquide et d'aucune option sur une action. Expliquer
ce qu'elle peut faire et quel risque subsiste.
\end{exercise}

\begin{exercise}[$\star\star\star$]\label{exo:m1:asian-and-emerging-equity-markets:7}
\emph{Programmation.} Avec \texttt{truncate}, générer 60\,000 rendements à queues lourdes
comme dans le tutoriel (graine 12), et calculer l'écart-type observé et
l'autocorrélation d'ordre un pour des limites de 5~\%, 10~\% et 20~\%. Expliquer cette régularité.
\end{exercise}

\begin{exercise}[$\star\star\star$]\label{exo:m1:asian-and-emerging-equity-markets:8}
\emph{Trouver la faille.} Un chercheur affirme~: «~Les A-shares qui clôturent \omterm{def:m1:asian-and-emerging-equity-markets:limit}{en limite haute}
gagnent en moyenne 3.1~\% de plus le lendemain~; acheter au cours de clôture \omterm{def:m1:asian-and-emerging-equity-markets:limit}{en limite haute}
rapporte 3.1~\% par jour.~» Identifier les deux raisons pour lesquelles ce profit n'est pas réalisable,
et la statistique qui aurait dû être rapportée.
\end{exercise}

\section{Problème~: Limit-Up}

\begin{problem}[{Problème du week-end --- une OPA sur un marché soumis à des limites}]\label{pb:m1:asian-and-emerging-equity-markets:1}
Une action clôture à 20.00 le lundi sur un marché doté d'une limite quotidienne de
variation de cours de 10~\% (les cours limites sont arrondis au centime inférieur), d'un \omterm{def:m1:asian-and-emerging-equity-markets:lot}{lot de négociation} de 100 actions, d'une interdiction
de vendre des actions achetées le jour même, et d'un droit de 0.1~\% sur chaque côté. Après
la clôture, un acquéreur annonce une offre en numéraire à 29.00, que le marché
juge certaine de se réaliser dans trois mois. Le taux d'intérêt est de 4~\% par
an.

\medskip\noindent\textbf{Partie I --- Le parcours.}
\begin{enumerate}
\item Quelle est la valeur de l'action le mardi matin, en actualisant l'offre au
taux d'intérêt pour trois mois~?
\item Donner le cours plafond le mardi, le mercredi, le jeudi et le vendredi,
chacun calculé à partir de la clôture précédente.
\item Quel est le premier jour où l'action peut se négocier à sa fair value~? Combien de jours
est-elle bloquée \omterm{def:m1:asian-and-emerging-equity-markets:limit}{en limite haute}~?
\item Confronter le résultat à la \cref{prop:m1:asian-and-emerging-equity-markets:locked}.
\end{enumerate}

\medskip\noindent\textbf{Partie II --- La file d'attente.}
Le mardi, des ordres d'achat portant sur 40~millions d'actions se placent en file au cours limite~; 300\,000
actions sont vendues à ce cours dans la journée, réparties selon la priorité temporelle. Un investisseur saisit
un ordre de 50\,000 actions 0.2~seconde après l'ouverture, derrière 9~millions d'actions
déjà en file.
\begin{enumerate}[resume]
\item Combien d'actions l'investisseur reçoit-il~?
\item Quelle est la valeur d'une action achetée au cours limite du mardi, par rapport à son
cours, en pourcentage~?
\item Un investisseur placé tout en tête de file achète 100\,000 actions. Quel est
le gain espéré en valeur monétaire, avant coûts~?
\item Qu'est-ce qui détermine qui se trouve en tête de file, et combien une société de trading
paierait-elle pour cela~?
\item S'il achète le mardi, quand peut-il vendre pour la première fois, et à quel cours limite
au mieux~?
\end{enumerate}

\medskip\noindent\textbf{Partie III --- Les données.}
\begin{enumerate}[resume]
\item Écrire la série des rendements quotidiens clôture-à-clôture, du lundi au
premier jour où l'action atteint la fair value, en supposant qu'elle clôture ensuite à la fair value.
\item Donner la somme des carrés des rendements quotidiens, et le carré du
rendement total~; comparer.
\item Un système de risque estime la volatilité de cette action à partir des 20 derniers
rendements quotidiens, les 16 autres étant nuls. Quelle volatilité quotidienne
le système affiche-t-il, et quelle serait-elle si le mouvement s'était produit en un seul jour~?
\item Calculer la première autocorrélation de votre série de la question 10
(sans retrancher la moyenne). Que conclurait un signal de momentum~?
\end{enumerate}

\medskip\noindent\textbf{Partie IV --- L'opération réalisable.}
\begin{enumerate}[resume]
\item Une fois l'action parvenue à la fair value, elle se négocie à 28.60 et une offre
à 29.00 à échéance dans trois mois. Donner le rendement brut de l'achat suivi de l'apport à l'offre,
annualisé.
\item Retrancher le \omterm{def:m1:asian-and-emerging-equity-markets:stamp}{droit de timbre} sur l'achat (l'apport étant une vente) et le
financement à 4~\%. L'opération en vaut-elle la peine~?
\item Le \omterm{def:m1:asian-and-emerging-equity-markets:lot}{lot de négociation} est de 100 actions et l'on souhaite investir 2 millions. Combien
d'actions cela représente-t-il~?
\item La même société dispose d'une ligne cotée sur un autre marché, sans \omterm{def:m1:asian-and-emerging-equity-markets:limit}{limite de variation
de cours}, et non convertible en celle-ci. Que faut-il attendre de cette ligne
sur la séance du mardi, et que mesure l'écart entre les deux lignes~?
\item Nommer le seul risque que l'hypothèse de réalisation certaine de l'offre a écarté, et
dire comment cela modifierait la question 14.
\item Énoncer le \emph{résultat nommé}~: le nombre de jours pendant lesquels l'action est bloquée
au cours limite.
\item En une phrase, que protège une \omterm{def:m1:asian-and-emerging-equity-markets:limit}{limite de variation de cours}, et
quel en est le coût~?
\end{enumerate}
\end{problem}

\section{Questions d'entretien}

\begin{interviewq}[$\star$ \iqroles{trader, chercheur}]\label{iq:m1:asian-and-emerging-equity-markets:1}
Qu'est-ce qu'une \omterm{def:m1:asian-and-emerging-equity-markets:limit}{limite de variation de cours} quotidienne, et que se passe-t-il pour une action lorsque
sa fair value varie de plus que la limite~?
\end{interviewq}

\begin{interviewq}[$\star$ \iqroles{chercheur, ingénieur ML}]\label{iq:m1:asian-and-emerging-equity-markets:2}
On vous fournit les cours de clôture quotidiens d'un marché doté de limites de variation de 10~\%. Citez deux
statistiques qui seraient biaisées si vous les traitiez comme des données américaines.
\end{interviewq}

\begin{interviewq}[$\star\star$ \iqroles{trader, chercheur}]\label{iq:m1:asian-and-emerging-equity-markets:3}
Un marché prélève 10 points de base de taxe sur chaque côté d'une transaction sur actions. Quelles
stratégies disparaissent, et qui finit par devenir fournisseur de liquidité~?
\end{interviewq}

\begin{interviewq}[$\star\star$ \iqroles{chercheur, trader}]\label{iq:m1:asian-and-emerging-equity-markets:4}
Les actions onshore et offshore d'une même société se négocient avec un écart de 30~\%. Pourquoi
cela n'est-il pas un arbitrage, et comment pourriez-vous néanmoins l'exploiter~?
\end{interviewq}

\begin{interviewq}[$\star\star$ \iqroles{développeur}]\label{iq:m1:asian-and-emerging-equity-markets:5}
Votre couche de validation des ordres doit prendre en charge quinze marchés dont les règles
de lot, de tick et de \omterm{def:m1:asian-and-emerging-equity-markets:limit}{limite de variation de cours} changent plusieurs fois par an. Comment la concevez-vous~?
\end{interviewq}

\begin{interviewq}[$\star\star\star$ \iqroles{chercheur, trader}]\label{iq:m1:asian-and-emerging-equity-markets:6}
Un régulateur interdit la vente à découvert pendant la nuit. Énumérez les stratégies et les
fonctions de marché touchées, par ordre de gravité, et que feriez-vous le premier
matin si vous gériez un portefeuille long--short là-bas~?
\end{interviewq}
